\section{Antecedentes y marco conceptual}

\subsection{Residuos orgánicos y bioconversión con \textit{Hermetia illucens}}

En contextos agroindustriales del sur global ---como el Valle del Cauca---, los residuos orgánicos provenientes de cultivos como café, caña y frutas se generan de forma dispersa, en picos estacionales y con alta carga biodegradable, lo que incrementa el riesgo de emisiones de \ce{CH4} cuando su gestión es inadecuada \cite{lacknerAgriculturalWasteChallenges2025, leddinImpactClimateChange2024a}. Frente a las limitaciones de las estrategias convencionales ---rellenos sanitarios, digestión anaerobia y compostaje---, que suelen ser inviables técnica o económicamente en estos entornos \cite{kongEvaluatingGreenhouseGas2012, nordahlGreenhouseGasAir2023a}, la bioconversión con BSF se ha adoptado crecientemente como una alternativa de economía circular.

La bioconversión de residuos orgánicos mediante larvas en estadios voraces (L3--L5) transforma desechos en biomasa rica en proteínas y lípidos, así como en frass, un fertilizante orgánico de alto valor agronómico \cite{bermudez-serranoComprehensiveUtilizationBlack2023, singhInclusiveApproachOrganic2019a, surendraBioconversionOrganicWastes2016}. Este proceso se sustenta en una sinergia entre la acción macroscópica de las larvas y la actividad microbiana del sustrato: las larvas fragmentan mecánicamente el material, aumentando su área superficial, mientras su microbioma intestinal ---junto con microorganismos del medio--- degrada enzimáticamente polímeros complejos en compuestos asimilables. Esta doble acción reduce entre el 50\% y el 80\% la masa del residuo (base húmeda) y altera drásticamente las propiedades fisicoquímicas del medio (pH, humedad, disponibilidad de nutrientes), generando microambientes dinámicos que condicionan si predominan rutas metabólicas aeróbicas o anaeróbicas. Durante este proceso, el dióxido de carbono (\ce{CO2}) se genera principalmente por la respiración aeróbica tanto de las larvas como de la microbiota asociada; estudios estiman que hasta un tercio del carbono ingerido se libera como \ce{CO2} \cite{dienerConversionOrganicMaterial2011}, una emisión altamente sensible a factores como la composición del sustrato, densidad larval, temperatura y oxigenación \cite{bekkerImpactSubstrateMoisture2021a, eriksenDynamicModellingFeed2022, rossiEstimatingDynamicsGreenhouse2024a}.

Aunque la bioconversión muestra un perfil de emisiones potencialmente más favorable que el compostaje tradicional ---con reducciones reportadas en \ce{CH4} y \ce{N2O} \cite{jenkinsProcessingPoultryManure2025}---, su implementación a escala carece de sistemas rigurosos para cuantificar y predecir en tiempo real las emisiones gaseosas asociadas. Esta limitación impide su validación como solución climática verificable y obstaculiza su integración en mecanismos de compensación de carbono o políticas de gestión sostenible de residuos. En particular, la dinámica de emisiones, especialmente la de \ce{CO2} ---que refleja directamente la actividad biológica en curso---, no ha sido suficientemente modelada mediante herramientas matemáticas capaces de ofrecer estimaciones continuas y predictivas bajo condiciones operativas variables. Así, persiste una brecha crítica entre el potencial ambiental demostrado en ensayos controlados y la capacidad de monitoreo, verificación y escalabilidad requerida para su reconocimiento en estrategias climáticas basadas en datos.

Aunque el metano (\ce{CH4}) es un gas de efecto invernadero con un potencial de calentamiento global significativamente mayor que el del dióxido de carbono (\ce{CO2}), su generación en sistemas de bioconversión con BSF es generalmente baja debido al carácter predominantemente aeróbico del proceso. No obstante, se incluye su monitoreo como variable de control para detectar condiciones transitoriamente anaeróbicas ---por ejemplo, por sobresaturación del sustrato o deficiente aireación--- que podrían comprometer la eficiencia del proceso y generar picos no deseados de emisiones. Sin embargo, el foco técnico de esta tesis recae en el \ce{CO2}, ya que constituye el principal indicador cuantitativo y dinámico de la actividad biológica total (tanto larval como microbiana) durante la bioconversión. Además, presenta una señal más estable y continua, lo que facilita su integración en modelos de estimación en tiempo real, y permite una correlación directa con variables operativas clave, como el consumo de sustrato, la biomasa generada y las tasas respiratorias. Por estas razones, el \ce{CO2} se adopta como la variable central para el desarrollo, calibración y validación del sistema híbrido propuesto.

La aptitud de los residuos para la bioconversión varía significativamente según su composición. Subproductos agroindustriales como la pulpa de café o cáscaras de frutas son altamente adecuados, mientras que residuos lignocelulósicos o lodos orgánicos requieren pretratamiento o presentan riesgos sanitarios \cite{brunoValorizationOrganicWaste2025, elsayedConversionProteinrichWaste2024a}. Esta heterogeneidad ---común en el Valle del Cauca--- subraya la necesidad de sistemas de monitoreo adaptables, capaces de caracterizar el proceso en función del sustrato y las condiciones locales.

La evidencia cuantitativa disponible es fragmentada: se reportan emisiones de \ce{CO2} entre 7.76 y 11.88 g por gramo de larva seca \cite{rossiEstimatingDynamicsGreenhouse2024a}, o 96.5 g de \ce{CO2} por kg de residuo tratado \cite{fuhrmannComprehensiveIndustryrelevantBlack2025a}, pero sin protocolos estandarizados ni mediciones continuas. Esta incertidumbre impide comparar tecnologías, diseñar políticas públicas informadas o acceder a mecanismos de mercado basados en la mitigación climática.

\subsection{Modelación de emisiones y sistemas inteligentes}

La cuantificación dinámica de emisiones en procesos biológicos exige herramientas que articulen una base mecanicista con capacidad de adaptación empírica. En este contexto, el modelado matemático basado en ecuaciones diferenciales ordinarias (EDO) ha impulsado avances significativos. Por ejemplo, el modelo de crecimiento larval propuesto por \citeauthor{eriksenDynamicModellingFeed2022} \cite{eriksenDynamicModellingFeed2022} integra balances energéticos y cinética logística para simular simultáneamente la acumulación de biomasa y la producción de \ce{CO2}, combinando los principios del enfoque Dynamic Energy Budget (DEB) con eventos biológicos clave del ciclo larval, como las mudas y la metamorfosis. De manera análoga, funciones cinéticas de tipo Monod y Haldane ---originalmente desarrolladas para digestión anaerobia--- han sido adaptadas con éxito para describir la degradación del sustrato y el crecimiento microbiano en sistemas de bioconversión con insectos \cite{ahlamineMathematicalAnalysisAnaerobic2024, winAnaerobicDigestionBlack2018a}.

Si bien los modelos mecanicistas actuales proporcionan una base teórica indispensable, su aplicación directa en entornos operativos dinámicos puede presentar limitaciones de alcance. La literatura \cite{bekkerImpactSubstrateMoisture2021a, rossiEstimatingDynamicsGreenhouse2024a} sugiere que, al validarse principalmente bajo condiciones controladas, las formulaciones clásicas a menudo requieren simplificaciones sobre las interacciones simultáneas de variables (temperatura, humedad, oxigenación). Asimismo, el uso de parámetros fijos podría restringir la capacidad del modelo para responder a la variabilidad estocástica del sustrato en campo. Esta diferencia entre las condiciones controladas y la operación real sugiere la conveniencia de explorar enfoques híbridos o basados en datos que complementen la teoría existente con información en tiempo real.

Paralelamente, los avances en sensores de bajo costo y plataformas IoT han facilitado el monitoreo continuo de variables críticas ---\ce{CO2}, \ce{CH4}, temperatura, humedad--- en entornos operativos reales \cite{duobieneDevelopmentWirelessSensor2022, rajRealTimeEstimation2023a, yavariArtEMonArtificialIntelligence2023}. Estos datos, procesados mediante algoritmos de aprendizaje automático (AAA), permiten identificar patrones complejos y no lineales que serían difíciles de capturar mediante reglas explícitas \cite{chenOptimizationModelProcess2021a, weichertReviewMachineLearning2019}.

En particular, la integración de modelos mecanicistas (EDO) con AAA, enfoque conocido en la literatura como modelado de caja gris (\textit{grey box modelling}), ha demostrado potencial en sistemas biotecnológicos complejos: mientras el modelo EDO proporciona una estructura interpretable y basada en principios físicos, el AAA permite calibración dinámica de parámetros a partir de datos en tiempo real, mejorando predicciones cuando las condiciones operativas se desvían del escenario de entrenamiento \cite{guillaumeAsymptoticEstimatedDigestibility2023, kobelskiProcessOptimizationBlack2024}.

Sin embargo, esta integración enfrenta desafíos prácticos: (i) los algoritmos de AAA requieren suficientes datos históricos para entrenar sin riesgo de \textit{overfitting}, pero ciclos de bioconversión típicos generan pocas muestras por evento; (ii) la validación cruzada y selección de hiperparámetros requieren cuidado especial en contextos con datos limitados; (iii) la elección entre arquitecturas simples (regresión lineal múltiple) versus complejas (redes neuronales) dependerá del volumen y calidad de datos disponibles. Por ello, esta tesis adopta un enfoque escalonado: comenzar con modelos interpretables, evaluar redes neuronales multicapa si se dispone de base de datos amplia, y validar todas las predicciones contra mediciones de referencia de campo.

\subsection{Gobernanza de datos y ética algorítmica}

La implementación de tecnologías digitales en la gestión de procesos biológicos no es un acto neutral; constituye la creación de un sistema socio-técnico donde convergen biología, código y decisión humana. Desde la perspectiva de los estudios de ciencia, tecnología y sociedad (CTS) y la epistemología ambiental, es imperativo trascender la visión instrumental de la tecnología para abordar las implicaciones éticas y cognitivas de ``datificar'' la naturaleza.

Uno de los mayores riesgos en la aplicación de la inteligencia artificial (IA) al medio ambiente es la opacidad algorítmica \cite{burrellHowMachineThinks2016}. En modelos de caja negra, las decisiones se vuelven inescrutables, generando relaciones de dependencia tecnológica donde el operario obedece sin comprender. Esta tesis adopta una postura ética fundada en la interpretabilidad por diseño. Al integrar modelos mecanicistas basados en ecuaciones diferenciales ordinarias (EDO) ---cuyos parámetros tienen significado biológico explícito--- con algoritmos de aprendizaje ligeros, el sistema no solo predice eventos (el qué), sino que permite rastrear sus causas fisiológicas o operativas (el porqué), como el exceso de humedad o la acumulación de zonas anaeróbicas. Esta transparencia cognitiva devuelve la agencia al tomador de decisiones y se alinea con los principios de transparencia exigidos por marcos de reporte climático como el IPCC Tier 3 \cite{adadiPeekingBlackBoxSurvey2018}.

La gobernanza de los datos ambientales plantea interrogantes críticos sobre propiedad, acceso y beneficio. La literatura sobre Big Data en agricultura alerta sobre el ``extractivismo de datos'', donde la información generada en territorios rurales es capturada por plataformas externas sin retorno de valor local \cite{carolanAutomatedAgrifoodFutures2020}. Esta investigación se alinea con los principios de soberanía de datos y diseño frugal. La arquitectura propuesta es descentralizada: el procesamiento crítico ocurre en el borde (\textit{edge computing}), minimizando la dependencia de infraestructuras en la nube centralizadas y costosas. Esto no solo reduce costos, sino que garantiza que los datos sensibles permanezcan bajo el control del productor.

Finalmente, se reconoce la tensión ontológica entre la complejidad del sistema biológico y su representación digital. Un sensor de \ce{CO2} registra una señal fisicoquímica, pero no captura dimensiones cualitativas como la resiliencia del microbioma o el bienestar larval. Existe el riesgo de caer en el solucionismo tecnológico, asumiendo que lo no medido es irrelevante. Por ello, este marco teórico establece que el sistema híbrido es una herramienta de soporte, no de sustitución. La comprensión más robusta del proceso emerge de la interacción dialógica entre la señal del sensor y la experiencia empírica del operador.

\subsection{Brechas tecnológicas y oportunidades de investigación}

A pesar de los avances, la implementación operativa de sistemas inteligentes en la bioconversión con BSF enfrenta barreras estructurales que limitan su escalabilidad y confiabilidad. La literatura identifica cinco brechas críticas que esta investigación busca abordar:

\begin{enumerate}
    \item \textbf{Materialidad y Resiliencia del Hardware:} Escasa disponibilidad de sensores de bajo costo capaces de mantener la precisión (calibración estable) en las condiciones agresivas del proceso (HR $>80\%$), lo que suele derivar en datos ruidosos o deriva instrumental \cite{AgricultureClimateChangea}.
    \item \textbf{Interoperabilidad:} Falta de estandarización en protocolos de comunicación que impide la integración fluida entre dispositivos IoT heterogéneos y plataformas de gestión ambiental \cite{vanIntegrationInternetofThingsSustainable2022a}.
    \item \textbf{Limitaciones de los Modelos de Caja Negra:} Los algoritmos de aprendizaje automático puros, aunque potentes, carecen de restricciones físicas o biológicas explícitas. Esto los hace propensos al \textit{overfitting} y reduce su capacidad de generalización ante condiciones no vistas (como cambios de sustrato), a menos que se hibriden con modelos mecanicistas (EDO) \cite{weichertReviewMachineLearning2019}.
    \item \textbf{Ausencia de Estándares MRV Dinámicos:} Inexistencia de metodologías validadas para la Medición, Reporte y Verificación (MRV) en tiempo real, lo que excluye a la BSF de los mercados de carbono regulados \cite{herrinCollateralDataQualitya}.
    \item \textbf{Brecha de Implementación Rural:} Dificultades para escalar soluciones dependientes de la nube en contextos con baja infraestructura de conectividad, exigiendo arquitecturas de procesamiento en el borde.
\end{enumerate}

Estas brechas delinean una oportunidad estratégica para el desarrollo de una arquitectura híbrida compuesta por: (i) una red de sensores robustecida para monitoreo ambiental; (ii) modelos EDO para capturar la dinámica biológica base; (iii) algoritmos de corrección por AAA para adaptarse a la variabilidad real; y (iv) una interfaz de soporte a la decisión. Esta integración apunta a transformar la bioconversión de una caja negra biológica a un proceso transparente, tecnológicamente viable y verificable climáticamente \cite{kamilarisReviewPracticeBig2017}.
